\begin{table}[H]
\centering
\caption{Output Token Analysis by Prompt Variant}
\label{tab:token_analysis}
\footnotesize
\begin{tabular}{@{}llrrrrrr@{}}
\toprule
\textbf{Variant} & \textbf{Description} & \textbf{Max Tok} & \textbf{Limit} & \textbf{Messages} & \textbf{Calls} & \textbf{Success} & \textbf{Rate} \\
\midrule
baseline & Direct evaluation & 1,836 & 1,500 & 27,649 & 34,660 & 34,641 & 99.9\% \\
mechanistic & Bradford Hill criteria & 291 & 1,500 & 9,269 & 9,262 & 9,262 & 100.0\% \\
mech_lit & Mechanistic (citation, literature) & 1,777 & 1,500 & 13,466 & 64,993 & 57,863 & 89.0\% \\
mech_orig & Mechanistic (citation, original) & 1,435 & 1,500 & 18,330 & 21,026 & 21,008 & 99.9\% \\
cot & Chain-of-thought reasoning & 1,856 & 10,000 & 24,273 & 59,795 & 59,304 & 99.2\% \\
\midrule
\textbf{Total} & --- & --- & --- & \textbf{92,987} & \textbf{189,736} & \textbf{182,078} & --- \\
\bottomrule
\end{tabular}
\begin{tablenotes}
\small
\item \textit{Note.} Max Tok = maximum observed output tokens (via tiktoken). Limit = max\_tokens parameter.
\item Calls = total API calls. Success = successful completions. Rate = success percentage.
\item \textbf{Conclusion:} All outputs below token limits. No truncation detected.
\end{tablenotes}
\end{table}